This work provides evidence for a Higgs mode that is stabilized by proximity to a quantum critical point. Dynamical measurements in CuBr2 detect a low-energy mode that is surprisingly long lived. The authors interpret this mode terms of a Higgs mode, which is made long-lived because a proximity to a quantum critical point makes the Higgs mode to lie below the Goldstone mode that is made to have a nonzero energy due to magnetic anisotropy. Theoretical speculations along this line have existed before (eg Ref 23).

Still, the experiment is striking and the mechanism brings out the general conceptual importance, viz new dynamical effect is enabled by quantum criticality. The fact that this general physics is happening in a class of materials that have not in the past been highlighted in the context of quantum criticality further elevates the importance. I therefore feel that the work is potentially suitable for publication in a high profile journal like Nature Materials.
